\section*{Bab \ref{ch:b1:body-plans-tissues} --- Pola Tubuh dan Jaringan Hewan}

\begin{solution}{exo:b1:body-plans-tissues:1}
Radial: tersusun mengelilingi sebuah sumbu, tanpa kiri dan kanan
(hidra, ubur-ubur). Bilateral: satu bidang cermin, sisi depan dan
belakang, punggung dan perut (cacing tanah, katak). \omterm{def:b1:mammal-organization:bodyplan}{Sefalisasi} sejalan
dengan \omterm{def:b1:body-plans-tissues:symmetry}{simetri bilateral}: hewan yang bergerak dengan kepala di depan
menemui dunia pada sisi depannya, tempat indra dan pusat saraf memusat.
\end{solution}

\begin{solution}{exo:b1:body-plans-tissues:2}
\omterm{def:b1:body-plans-tissues:layers}{Ektoderm}: \omterm{def:b1:flowering-plant-organization:tissues}{epidermis} kulit, \omterm{def:b1:body-plans-tissues:nervous}{jaringan saraf}. \omterm{def:b1:body-plans-tissues:layers}{Mesoderm}: otot, tulang,
darah, \omterm{def:b1:body-plans-tissues:connective}{jaringan ikat}. \omterm{def:b1:body-plans-tissues:layers}{Endoderm}: pelapis usus.
\end{solution}

\begin{solution}{exo:b1:body-plans-tissues:3}
\omterm{def:b1:body-plans-tissues:epithelium}{Epitel}: sel \omterm{def:b1:body-plans-tissues:epithelium}{terkutub} yang tersusun rapat sebagai lembar di atas \omterm{def:b1:body-plans-tissues:epithelium}{lamina
basal}; penyelimutan, penyerapan, sekresi. Ikat: sel yang tersebar dalam
matriks yang melimpah; pengikatan, penopangan, pengangkutan. Otot: sel
yang padat berisi filamen yang dapat berkerut; gerak. Saraf: \omterm{def:b1:body-plans-tissues:nervous}{neuron} dan
glia; penginderaan, pemaduan, perintah.
\end{solution}

\begin{solution}{exo:b1:body-plans-tissues:4}
Sambungan rapat, yang menutup celah antarsel; sambungan lekat, yang
mengikat dan terhubung ke sabuk aktin; desmosom, paku keling mekanis
yang terhubung ke filamen antara; sambungan celah, yaitu saluran di
antara kedua sitoplasmanya.
\end{solution}

\begin{solution}{exo:b1:body-plans-tissues:5}
Triploblastik, selomata, berusus tembus, \omterm{def:b1:body-plans-tissues:segments}{bersegmen}, bertubuh lunak
dengan \omterm{def:b1:body-plans-tissues:segments}{rangka hidrostatik}: sebuah anelida (cacing tanah).
\end{solution}

\begin{solution}{exo:b1:body-plans-tissues:6}
Cairan \omterm{def:b1:body-plans-tissues:layers}{selom} tiap segmen tidak dapat dimampatkan. Mengerutkan otot
melingkar membuat segmen itu kurus dan panjang (ia mendorong ke depan);
mengerutkan otot memanjang membuatnya pendek dan gemuk (ia menambat pada
dinding liang dengan rambut kejurnya). Gelombang kedua kerutan yang
berjalan mundur sepanjang tubuh menarik cacing itu maju.
\end{solution}

\begin{solution}{exo:b1:body-plans-tissues:7}
\omterm{def:b1:body-plans-tissues:segments}{Rangka luar}: kaku, melindungi, otot melekat di sisi dalam; tidak dapat
tumbuh, sehingga hewan itu harus berganti kulit dan sementara itu lunak
serta rentan; di darat massanya tumbuh sebagai pangkat tiga ukurannya
sedangkan kekuatannya sebagai pangkat dua, yang membatasi artropoda pada
ukuran kecil. \omterm{def:b1:body-plans-tissues:segments}{Rangka dalam}: tumbuh bersama hewannya, tanpa pergantian
kulit, perlindungan permukaan lebih sedikit, otot melekat di sisi luar;
memungkinkan ukuran yang besar.
\end{solution}

\begin{solution}{exo:b1:body-plans-tissues:8}
Sambungan rapat memaksa tiap zat melewati sel, yang membrannya membawa
pengangkut yang memilih; \omterm{def:b1:body-plans-tissues:epithelium}{epitel} itu lalu dapat memilih apa yang
melintas dan ke arah mana. Tanpa sambungan itu isi lumen akan merembes
di antara sel masuk ke \omterm{def:b1:organism-environment:milieu}{lingkungan internal} (bakteri, racun, molekul yang
belum tercerna), dan zat hara yang terserap akan merembes kembali.
\end{solution}

\begin{solution}{exo:b1:body-plans-tissues:9}
Usus halus. Mukosa (\omterm{def:b1:body-plans-tissues:epithelium}{epitel} dengan vili, \omterm{def:b1:body-plans-tissues:connective}{jaringan ikat} longgarnya,
muscularis mucosae yang tipis), submukosa (\omterm{def:b1:body-plans-tissues:connective}{jaringan ikat} padat),
muskularis (otot polos melingkar di dalam, memanjang di luar),
serosa.
\end{solution}

\begin{solution}{exo:b1:body-plans-tissues:10}
Sel tulang rawan diberi makan lewat difusi menembus gelnya dari
pembuluh di sekelilingnya, sehingga \omterm{def:b1:body-plans-tissues:tissue}{jaringan} itu tipis atau
bermetabolisme lambat, sukar pulih dan dipakai di tempat yang menuntut
permukaan licin, mampat dan tanpa pembuluh: permukaan sendi, telinga,
hidung, ujung tulang yang sedang tumbuh. Sel tulang duduk di dalam
mineral yang tak tertembus difusi apa pun dan harus diberi makan lewat
pembuluh di dalam kanal; tulang pulih dengan baik dan dibangun ulang
seumur hidup, serta membentuk rangka pemikul beban.
\end{solution}

\begin{solution}{exo:b1:body-plans-tissues:11}
Penyerapan menjadi tak memilih: zat hara melintas di antara sel
menuruni gradiennya dan kerja berarah pengangkutnya hilang, sementara
molekul yang telah terserap merembes kembali. \omterm{def:b1:organism-environment:milieu}{Lingkungan internal}
menerima apa pun yang ada di dalam lumen --- ion, air, peptida yang
belum tercerna, hasil bakteri --- dan kehilangan air serta ion ke dalam
usus: diare, udem atau dehidrasi, susunan plasma yang terganggu.
Bakteri dan racun yang melintasi sawar itu menimbulkan radang dan
infeksi.
\end{solution}

\begin{solution}{exo:b1:body-plans-tissues:12}
Sebuah \omterm{def:b1:mammal-organization:bodyplan}{pola tubuh} diwariskan, bukan dipilih menurut kebutuhan hewannya,
dan pola itu menetapkan apa yang dapat dilakukan evolusi kemudian:
\omterm{def:b1:body-plans-tissues:segments}{rangka luar} menuntut pergantian kulit dan menahan serangga tetap kecil;
seekor vertebrata harus membangun tubuhnya dari segmen bahkan ketika
segmen itu tak berguna baginya; hewan bersimetri radial sulit
mengembangkan sebuah kepala. Tabel tadi mendaftar pilihan yang
diserahkan kepada tiap garis keturunan, masing-masing lengkap dan
berhasil di dalam kendalanya --- bukan sebuah tangga, sebab pola katak
bukanlah perbaikan atas pola udang karang, melainkan hanya sehimpunan
kemungkinan yang lain.
\end{solution}

\begin{solution}{pb:b1:body-plans-tissues:1}
\textbf{1.} Mukosa: \omterm{def:b1:body-plans-tissues:epithelium}{epitel}, \omterm{def:b1:body-plans-tissues:connective}{jaringan ikat} longgar, selembar otot polos.
Submukosa: \omterm{def:b1:body-plans-tissues:connective}{jaringan ikat} padat (pembuluh, saraf). Muskularis: otot polos
dalam dua lapis, \omterm{def:b1:body-plans-tissues:nervous}{jaringan saraf} di antaranya. Serosa: \omterm{def:b1:body-plans-tissues:connective}{jaringan ikat}
dengan sebuah \omterm{def:b1:body-plans-tissues:epithelium}{epitel}.
\textbf{2.} Usus halus; mukosa lambung tidak memiliki vili melainkan
kelenjar tabung yang dalam.
\textbf{3.} Tubuh \omterm{def:b1:body-plans-tissues:nervous}{neuron} jala saraf usus itu sendiri; keduanya
mengendalikan kerutannya (pengadukan dan pendorongan) dan sekresinya.
\textbf{4.} Lapis melingkar menyempitkan tabungnya (\omterm{def:b1:body-plans-tissues:segments}{segmentasi},
pemerasan); lapis memanjang memendekkannya. Pergantian keduanya
mendorong isinya maju.
\textbf{5.} $\pi\times 0.03\times 6 = \qty{0.57}{m^2}$.
\textbf{6.} \qty{1.7}{m^2}.
\textbf{7.} $\pi\times 0.15\times 0.6 = \qty{0.28}{mm^2}$.
\textbf{8.} $25\times 0.28 = \qty{7.1}{mm^2}$ tiap mm$^2$: faktor $7$.
\textbf{9.} $1.7\times 7.1 = \qty{12}{m^2}$.
\textbf{10.} $\pi\times 0.1\times 1 = \qty{0.31}{\micro m^2}$;
\qty{470}{\micro m^2} tiap sel.
\textbf{11.} Sisi apikal \qty{25}{\micro m^2}: faktor $19$.
\textbf{12.} $12\times 19 = \qty{230}{m^2}$.
\textbf{13.} $\qty{0.28}{mm^2}/\qty{25}{\micro m^2} = \num{11000}$ sel.
\textbf{14.} $1.7\,\mathrm{m^2}\times 25\times 10^6\,\mathrm{mm^{-2}} =
\num{4.2e7}$ vilus; $\num{4.7e11}$ sel.
\textbf{15.} $\num{1.2e11}$ tiap hari, $\num{1.4e6}$ tiap sekon.
\textbf{16.} Di dalam kelenjar (kripta) di pangkal vilinya; sel yang
membelah terlindung di dasarnya dan hasilnya bermigrasi naik sepanjang
vilus, menyerap sambil berjalan, lalu dilepas dari ujungnya.
\textbf{17.} $12\,\mathrm{m^2}\times 25\times 10^{-6}\,\mathrm{m} =
\qty{3e-4}{m^3} = \qty{300}{cm^3}$, sekitar \qty{315}{g}.
\textbf{18.} Kelilingnya \qty{20}{\micro m}, yakni \qty{10}{\micro m}
tiap sel (tiap tepi dipakai bersama);
$\num{4.7e11}\times 10^{-5} = \qty{4700}{km}$.
\textbf{19.} Sambungan rapatnya menutup lintasan antarsel; zat itu harus
masuk ke sebuah sel lewat membran apikalnya dan keluar lewat membran
basalnya.
\textbf{20.} \omterm{def:b1:body-plans-tissues:connective}{Kolagen} dan \omterm{def:b1:body-plans-tissues:connective}{elastin}, yang disekresikan fibroblas.
\textbf{21.} \omterm{def:b1:body-plans-tissues:epithelium}{Epitel}: \omterm{def:b1:body-plans-tissues:layers}{endoderm}; teras (\omterm{def:b1:body-plans-tissues:connective}{jaringan ikat}, pembuluh, otot):
\omterm{def:b1:body-plans-tissues:layers}{mesoderm}.
\textbf{22.} Pemompaan: memendekkan vilus memeras pembuluh limfa dan
kapilernya serta memperbarui cairan di terasnya, sehingga gradien
melintasi \omterm{def:b1:body-plans-tissues:epithelium}{epitel} tetap curam.
\textbf{23.} Sebesar faktor vilusnya, $7$ (mikrovili \omterm{def:b1:body-plans-tissues:epithelium}{epitel} yang
merata tetap ada): salah serap lemak, gula dan vitamin, turun berat
badan, diare.
\textbf{24.} $0.28\times 19 = \qty{5.4}{mm^2}$.
\textbf{25.} Sekitar \qty{5}{mm^2} membran penyerap pada sebuah vilus
setinggi \qty{0.6}{mm}, dan sekitar empat puluh juta vilus dalam usus
halus.
\end{solution}
